%!TEX root = ../main.tex
\section{Идея 1: Рекурсивное деление пути (Recursive Bisection)}

\subsection{Как я это придумал и почему захотел попробовать сразу}

Первая мысль, которая пришла мне в голову при переходе к последовательностям — это бисекция пути пополам как в быстрой сортировке будем рекурсивно делить. Зачем что-то выдумывать и обучать  если можно в лоб найти в оффлайн-буфере такое промежуточное состояние $w^*$, до которого максимально легко дойти из текущей позиции $s$, и из которого сама цель $g$ также максимально достижима?  Идея проста и легко реализуема тч с неё я и решил начать.

Формально это сводится к максимизации логарифмической суммы достижимостей по пулу ориентиров:
\addequation{
    w^* = \arg\max_{w} \left[ \log(F(s, B(w))^\top B(w)) + \log(F(w, B(g))^\top B(g)) \right].
}{eq:recursive_bisection_objective}

Найдя точку $w^*$, мы рекурсивно делим отрезки $(s, w^*)$ и $(w^*, g)$ до глубины $D=2..3$, получая цепочку подцелей $[w_1, w_2, \dots, g]$ (алгоритм~\ref{alg:recursive_bisection}).

\begin{algorithm}[t]
\SetAlgoLined
\KwIn{Состояние $s$, цель $B(g)$, глубина $D$, порог $\tau$}
\KwOut{Цепочка подцелей $\mathcal{P}$}
\BlankLine
\SetKwFunction{FRec}{Bisect}
\SetKwProg{Fn}{Функция}{:}{}
\Fn{\FRec{$s_a$, $s_b$, $d$}}{
    \If{$d = 0$ \KwOr $F(s_a, B(s_b))^\top B(s_b) > \tau$}{
        \Return $[s_b]$\;
    }
    $w^* \leftarrow \arg\max_{w} [ \log F(s_a, B(w))^\top B(w) + \log F(w, B(s_b))^\top B(s_b) ]$\;
    \Return \FRec{$s_a$, $w^*$, $d-1$} + \FRec{$w^*$, $s_b$, $d-1$}\;
}
$\mathcal{P} \leftarrow$ \FRec{$s$, $g$, $D$}\;
\caption{Рекурсивная бисекция}
\label{alg:recursive_bisection}
\end{algorithm}

В процессе тестов выяснилось, что работает он хуже бейзлайна.  Набрав всего $75.9 \pm 6.4$ \% SR. И вот моё объяснение этому.

\subsection{Почему метод провалился: наступаем на первые грабли}

В главная проблема этого алгоритма в том, что он лишь усугубляет проблему отсутсвия локального контекста у планировщика. Например наиболее достяжимой точкой в начале окажется точка в U образной форме илич то-то в этом роде. проделав так ещё несколько раз мы получим троекторию которая очень вероятно будет идти через стены. И если взглянуть на характерные картинки то мы вообще увидем, что выходит не троектория а кластер точек тч именно они оказываються максимально достяжимы с двух краёв. Так что я сделал вывод, что в таком виде алгоритм несостоятелен и оставил его в покое.


\begin{figure}[t]
    \centering
    \begin{tikzpicture}[>=stealth, font=\small, scale=0.95, every node/.style={transform shape}]
    % Лабиринт с U-образным тупиком
    \draw[thick, fill=gray!10] (0,0) rectangle (7,4);
    
    % U-образная стена
    \draw[ultra thick, line width=3.5pt, red!70!black] (2.2, 0.8) -- (2.2, 3.2) -- (4.8, 3.2) -- (4.8, 0.8);
    \node[font=\footnotesize\bfseries, red!70!black] at (3.5, 3.45) {U-образное препятствие};
    
    % Точки
    \node[circle, fill=green!70!black, inner sep=2.5pt, label=below:{\footnotesize Старт $s$}] (S) at (1.2, 1.5) {};
    \node[star, star points=5, star point ratio=2.25, fill=yellow!80!black, draw=black, inner sep=2.5pt, label=below:{\footnotesize Финиш $g$}] (G) at (5.8, 1.5) {};
    
    % Ошибочный выбор точки деления пополам
    \node[circle, fill=orange!90!black, inner sep=2.5pt, label=above:{\footnotesize Ошибочный $w^*$}] (W) at (3.5, 2.2) {};
    
    \draw[->, thick, dashed, orange!80!black] (S) -- (W);
    \draw[->, thick, dashed, orange!80!black] (W) -- (G);
    \node[cross out, draw=red, line width=2pt, minimum size=10pt] at (2.2, 1.85) {};
    
    % Истинный маршрут обхода
    \draw[->, ultra thick, blue!80!black] (S) -- (1.2, 0.4) -- (5.8, 0.4) -- (G);
    \node[fill=white, inner sep=2pt, font=\scriptsize, align=center, blue!80!black] at (3.5, 0.15) {Истинный маршрут обхода снизу};

\end{tikzpicture}

    \caption{Типичный капкан бисекции: в U-образном препятствии точка $w^*$ выбирается прямо в середине тупика за стеной.}
    \label{fig:bisection_failure}
\end{figure}
